\documentclass[12pt,a4paper]{article}

\usepackage[T1]{fontenc}
\usepackage[utf8]{inputenc}
\usepackage[french]{babel}
\usepackage{geometry}
\usepackage{hyperref}
\usepackage{booktabs}
\usepackage{longtable}
\usepackage{tabularx}
\usepackage{enumitem}

\geometry{margin=2.5cm}
\hypersetup{
  colorlinks=true,
  linkcolor=black,
  urlcolor=blue,
  citecolor=black
}

\title{Cahier des charges \\ \large FrenchVerbs --- Cartes éducatives de conjugaison française}
\author{}
\date{\today}

\begin{document}
\maketitle
\tableofcontents
\newpage

\section{Contexte et objectifs}
Le projet \textbf{FrenchVerbs} est une application web destinée à faciliter l'apprentissage de la conjugaison française grâce à un système de \textbf{cartes éducatives} inspirées du style \emph{rami}. L'utilisateur authentifié peut consulter, filtrer et rechercher des verbes, accéder à une fiche détaillée, imprimer des cartes, gérer ses favoris et créer de nouvelles cartes. L'administration dispose d'outils pour \textbf{personnaliser le thème} visuel des cartes et superviser l'état du système.

\subsection{Objectifs pédagogiques}
\begin{itemize}[leftmargin=*]
  \item Proposer un support visuel clair et mémorisable pour apprendre le présent de l'indicatif.
  \item Permettre l'entraînement par répétition via une navigation rapide entre cartes.
  \item Offrir des fonctionnalités d'accessibilité (clavier, retours visuels) et d'assistance (guide débutant).
\end{itemize}

\subsection{Objectifs techniques}
\begin{itemize}[leftmargin=*]
  \item Application Laravel (PHP) avec pages Blade et interactions JavaScript.
  \item API JSON permettant d'exposer les verbes et leur conjugaison au présent.
  \item Zone d'administration sécurisée pour personnaliser et maintenir l'application.
\end{itemize}

\section{Périmètre}

\subsection{Inclus}
\begin{itemize}[leftmargin=*]
  \item Gestion des cartes de verbes au présent (liste, recherche, filtre, détail, impression, impression deck).
  \item Création d'une carte avec validation côté serveur et outils d'édition (aperçu live, modèles, annuler/rétablir, tutoriel, brouillon).
  \item Duplication d'une carte existante pour accélérer la création.
  \item Favoris : ajout/retrait et liste dédiée.
  \item Quiz : page d'entraînement (sélection de groupes).
  \item Authentification (inscription, connexion, déconnexion) et rôle administrateur.
  \item Administration : page Système (statut, actions) et personnalisation du thème (variables CSS).
  \item API JSON: récupération de la liste des verbes et d'un verbe.
  \item Expérience utilisateur : notifications toast, thème clair/sombre, anti double-soumission, TTS.
\end{itemize}

\subsection{Exclus / Hors périmètre}
\begin{itemize}[leftmargin=*]
  \item Gestion avancée des temps (imparfait, futur, etc.) au-delà du présent.
  \item Téléversement/gestion d'illustrations via interface (l'application supporte un chemin d'illustration mais pas de module d'upload dédié).
  \item Internationalisation complète (l'application est principalement en français).
\end{itemize}

\section{Acteurs et profils}
\begin{longtable}{@{}p{3.5cm}p{11cm}@{}}
\toprule
\textbf{Acteur} & \textbf{Description} \\
\midrule
\endhead
Visiteur & Utilisateur non authentifié : accès à la page d'accueil, inscription/connexion. \\
Utilisateur & Utilisateur authentifié : consultation des cartes, favoris, quiz, création/duplication et impression. \\
Administrateur & Utilisateur authentifié avec \texttt{is\_admin = true} : accès aux pages \texttt{/admin/theme} et \texttt{/admin/system}. \\
\bottomrule
\end{longtable}

\section{Architecture générale}

\subsection{Technologies}
\begin{itemize}[leftmargin=*]
  \item Backend : Laravel Framework (PHP), contrôleurs HTTP, validation serveur, Eloquent ORM.
  \item Frontend : Blade, CSS (fichier principal), JavaScript (jQuery) pour interactions.
  \item Outils frontend : Vite (build), Tailwind en dépendance (projet compatible Vite).
  \item Base de données : MySQL/PostgreSQL/SQLite selon configuration Laravel (migrations fournies).
\end{itemize}

\subsection{Pages et routes principales}
\begin{longtable}{@{}p{4cm}p{3cm}p{7.5cm}@{}}
\toprule
\textbf{Route} & \textbf{Méthode} & \textbf{Rôle} \\
\midrule
\endhead
\texttt{/} & GET & Page d'accueil (\emph{landing}). \\
\texttt{/dashboard} & GET & Redirection vers l'accueil (zone post-connexion). \\
\texttt{/cards} & GET & Liste des cartes (pagination, recherche \texttt{q}, filtre \texttt{group}). \\
\texttt{/cards/create} & GET & Formulaire de création d'une carte. \\
\texttt{/builder} & GET & Alias vers la création de carte. \\
\texttt{/cards} & POST & Création d'un verbe et de sa conjugaison au présent. \\
\texttt{/cards/\{verb\}/duplicate} & GET & Pré-remplissage du formulaire depuis un verbe existant. \\
\texttt{/cards/\{verb\}} & GET & Détail d'une carte (onglets pronoms). \\
\texttt{/cards/\{verb\}/print} & GET & Page d'impression (cartes par pronom). \\
\texttt{/cards/print/deck} & GET & Impression deck 9 cartes/A4 (options groupe, recto/verso). \\
\texttt{/rules} & GET & Règles et aide (page dédiée). \\
\texttt{/favorites} & GET & Liste des cartes favorites de l'utilisateur. \\
\texttt{/favorites/\{verb\}/toggle} & POST & Ajout/retrait d'un favori. \\
\texttt{/quiz} & GET & Page quiz (sélection de groupes). \\
\texttt{/login} & GET/POST & Connexion (throttle). \\
\texttt{/register} & GET/POST & Inscription (throttle). \\
\texttt{/logout} & POST & Déconnexion (auth). \\
\texttt{/health} & GET & Endpoint JSON de santé (DB + storage). \\
\texttt{/up} & GET & Endpoint de santé Laravel (configuré par l'application). \\
\texttt{/admin/verbs} & GET & Admin : liste des verbes et activation/désactivation. \\
\texttt{/admin/theme} & GET/POST & Admin : éditeur thème et sauvegarde des variables CSS. \\
\texttt{/admin/theme/reset} & POST & Admin : remise à zéro des valeurs de thème. \\
\texttt{/admin/system} & GET & Admin : tableau statut système. \\
\texttt{/admin/system/action} & POST & Admin : exécution d'actions (clear caches). \\
\bottomrule
\end{longtable}

\subsection{API JSON}
\begin{longtable}{@{}p{5cm}p{2cm}p{7.5cm}@{}}
\toprule
\textbf{Endpoint} & \textbf{Méthode} & \textbf{Description} \\
\midrule
\endhead
\texttt{/api/verbs} & GET & Liste des verbes actifs (filtres \texttt{group}, \texttt{q}). Retourne les données du verbe et la conjugaison au présent. \\
\texttt{/api/verbs/\{verb\}} & GET & Détails d'un verbe (inclut la conjugaison au présent formatée). \\
\bottomrule
\end{longtable}

\section{Exigences fonctionnelles}

\subsection{Module 1 --- Consultation des cartes}
\begin{itemize}[leftmargin=*]
  \item Afficher une liste paginée des verbes actifs (24 par page).
  \item Permettre la recherche par infinitif ou traduction via un champ \texttt{q}.
  \item Permettre le filtrage par groupe verbal \texttt{1er}, \texttt{2ème}, \texttt{3ème}.
  \item Afficher, pour chaque carte, un badge groupe et un symbole (\(\spadesuit\), \(\diamondsuit\), \(\clubsuit\), \(\heartsuit\) pour les auxiliaires \emph{être/avoir}).
  \item Permettre l'accès au détail d'une carte (navigation via clic/touche entrée).
\end{itemize}

\subsection{Module 2 --- Détail d'une carte}
\begin{itemize}[leftmargin=*]
  \item Afficher une carte \emph{grande} du verbe choisi et sa conjugaison au présent.
  \item Proposer des onglets de sélection du pronom (\texttt{JE, TU, IL, NOUS, VOUS, ILS}) via paramètre \texttt{p}.
  \item Afficher la description d'illustration si disponible.
  \item Proposer un bouton de prononciation (TTS) du verbe.
  \item Proposer un accès à l'impression de la carte.
\end{itemize}

\subsection{Module 3 --- Impression}
\begin{itemize}[leftmargin=*]
  \item Fournir une page d'impression générant une grille de cartes (une par pronom).
  \item Masquer l'interface non nécessaire (\emph{chrome}) en mode impression.
  \item Fournir des actions \emph{Retour} et \emph{Imprimer} (non imprimées).
  \item Fournir une impression \emph{deck} 9 cartes/A4 (mélange par groupes, limite pages/cartes, option verso).
\end{itemize}

\subsection{Module 4 --- Création d'une carte}
\begin{itemize}[leftmargin=*]
  \item Afficher un formulaire de création comprenant : infinitif, traduction, groupe, couleurs (thème/accent), motif, conjugaisons (je/tu/il-nous-vous-ils).
  \item Proposer un aperçu en temps réel de la carte en cours de saisie.
  \item Proposer une action \emph{Conjugaison Magique} pour suggérer automatiquement des formes (notamment 1er/2ème groupe) et informer l'utilisateur en cas d'impossibilité.
  \item Proposer des outils d'édition : modèles (presets), annuler/rétablir, tutoriel guidé, réinitialisation.
  \item Sauvegarder un brouillon côté navigateur (localStorage) et restaurer automatiquement si le formulaire est vide.
  \item Permettre la duplication d'une carte existante pour préremplir le formulaire.
  \item Appliquer des validations serveur :
    \begin{itemize}[leftmargin=*]
      \item Champs requis pour l'infinitif et toutes les formes au présent.
      \item Vérification du suffixe de l'infinitif selon le groupe (1er \(\rightarrow\) \texttt{er}, 2ème \(\rightarrow\) \texttt{ir}).
      \item Unicité de l'infinitif (insensible à la casse).
      \item Couleurs au format hexadécimal \texttt{\#RRGGBB}.
      \item Motif limité à une liste autorisée.
    \end{itemize}
  \item Enregistrer en base : un \emph{verb} et une \emph{conjugation} au temps \emph{présent}.
  \item Afficher un message de succès après création.
\end{itemize}
\subsection{Module 5 --- Favoris}
\begin{itemize}[leftmargin=*]
  \item Permettre d'ajouter/retirer un verbe des favoris (toggle).
  \item Afficher une page listant les favoris avec filtre groupe et recherche.
\end{itemize}

\subsection{Module 6 --- Quiz}
\begin{itemize}[leftmargin=*]
  \item Fournir une page quiz permettant de choisir un ou plusieurs groupes.
\end{itemize}

\subsection{Module 7 --- Authentification et sessions}
\begin{itemize}[leftmargin=*]
  \item Inscription avec \texttt{name}, \texttt{email}, \texttt{password} (confirmation).
  \item Connexion avec \texttt{email} et \texttt{password}, option \emph{remember}.
  \item Déconnexion via requête POST (CSRF).
  \item Limitation anti-bruteforce (throttle) sur les routes de connexion/inscription.
  \item Rôle administrateur : champ \texttt{is\_admin} sur l'utilisateur.
\end{itemize}

\subsection{Module 8 --- Administration : personnalisation du thème}
\begin{itemize}[leftmargin=*]
  \item Accès réservé aux administrateurs (\texttt{auth} + middleware \texttt{admin}).
  \item Interface de personnalisation basée sur des \textbf{variables CSS} (prévisualisation en direct).
  \item Sauvegarde via requête AJAX des valeurs autorisées, avec validation des formats.
  \item Réinitialisation du thème aux valeurs par défaut.
  \item Mise en cache des styles générés (CSS inline injecté dans le layout) et invalidation du cache à la sauvegarde.
  \item Préréglages : ombres (shadow presets), style d'illustration (center presets), icône d'illustration.
  \item Aide UX : recherche de paramètres, raccourcis clavier (ex: \texttt{Ctrl+S} pour enregistrer), avertissement avant de quitter si modifications non sauvegardées.
\end{itemize}

\subsection{Module 9 --- Administration : page Système}
\begin{itemize}[leftmargin=*]
  \item Afficher l'état général : nom app, environnement, debug, versions Laravel/PHP, drivers (DB/cache/queue/session), accès DB, droits d'écriture storage/logs.
  \item Proposer des actions d'entretien :
    \begin{itemize}[leftmargin=*]
      \item vider cache thème (cache applicatif),
      \item \texttt{optimize:clear}, \texttt{cache:clear}, \texttt{config:clear}, \texttt{route:clear}, \texttt{view:clear}.
    \end{itemize}
  \item Afficher les erreurs en message utilisateur en cas d'échec d'une action.
\end{itemize}

\subsection{Module 10 --- Monitoring}
\begin{itemize}[leftmargin=*]
  \item Endpoint \texttt{/health} renvoyant JSON :
    \begin{itemize}[leftmargin=*]
      \item \texttt{status} (\texttt{ok} ou \texttt{fail}),
      \item \texttt{timestamp},
      \item \texttt{checks} : base de données et écriture storage.
    \end{itemize}
  \item Code HTTP 200 si OK, 503 sinon.
\end{itemize}

\subsection{Module 11 --- Expérience utilisateur (UX) et accessibilité}
\begin{itemize}[leftmargin=*]
  \item Thème clair/sombre : bascule via bouton UI, persistance par \texttt{localStorage}.
  \item Notifications toast globales pour retours rapides (succès/erreur).
  \item Guide débutant (affichable/masquable, mémorisé en \texttt{localStorage}).
  \item Animations à l'apparition (\texttt{IntersectionObserver}) sur certaines cartes.
  \item Navigation clavier : cartes et éléments interactifs accessibles (\texttt{tabindex}, \texttt{role}).
  \item Anti double-soumission : blocage de formulaire au submit et désactivation des boutons.
  \item Copie facilitée : clic sur champs hex readonly et feedback toast.
\end{itemize}

\section{Exigences non fonctionnelles}
\subsection{Sécurité}
\begin{itemize}[leftmargin=*]
  \item Protection CSRF sur les formulaires.
  \item Validation serveur systématique des entrées utilisateur.
  \item Contrôle d'accès administrateur basé sur \texttt{is\_admin}.
  \item (Optionnel selon environnement) Basic Auth admin via variables d'environnement (middleware dédié, non obligatoire).
\end{itemize}

\subsection{Performance}
\begin{itemize}[leftmargin=*]
  \item Pagination sur la liste (limitation du volume de données).
  \item Cache des styles de thème (éviter calcul répété).
\end{itemize}

\subsection{Qualité et maintenabilité}
\begin{itemize}[leftmargin=*]
  \item Couverture par tests automatisés (ex: accès admin, validation création, endpoint health).
  \item Conventions Laravel (controllers, models, migrations, providers).
\end{itemize}

\subsection{Compatibilité}
\begin{itemize}[leftmargin=*]
  \item Fonctionnement sur navigateurs modernes (TTS et IntersectionObserver avec fallback implicite).
  \item Compatibilité DB multi-driver via migrations.
\end{itemize}

\section{Modèle de données (base)}
\begin{longtable}{@{}p{3.5cm}p{11cm}@{}}
\toprule
\textbf{Table} & \textbf{Champs principaux} \\
\midrule
\endhead
\texttt{users} & \texttt{name}, \texttt{email}, \texttt{password}, \texttt{is\_admin} \\
\texttt{verbs} & \texttt{infinitive}, \texttt{infinitive\_translation}, \texttt{group}, \texttt{illustration\_path}, \texttt{illustration\_description}, \texttt{theme\_color}, \texttt{accent\_color}, \texttt{pattern}, \texttt{is\_active} \\
\texttt{conjugations} & \texttt{verb\_id}, \texttt{tense}, \texttt{je}, \texttt{tu}, \texttt{il\_elle\_on}, \texttt{nous}, \texttt{vous}, \texttt{ils\_elles} \\
\texttt{theme\_settings} & \texttt{key} (unique), \texttt{value} \\
\texttt{cache}, \texttt{jobs} & Tables techniques Laravel (selon configuration). \\
\bottomrule
\end{longtable}

\section{Règles de gestion}
\begin{itemize}[leftmargin=*]
  \item Un verbe est affiché si \texttt{is\_active = true}.
  \item L'infinitif est unique (insensible à la casse).
  \item Le temps manipulé par défaut est \textbf{présent}.
  \item Un verbe appartient à un groupe \texttt{1er/2ème/3ème}.
  \item Les clés de personnalisation thème acceptées sont restreintes à une liste contrôlée (variables CSS \texttt{--rami-*}).
  \item Les valeurs de thème sont filtrées pour éviter l'injection (caractères \texttt{; \{ \}} interdits, longueurs limitées).
\end{itemize}

\section{Cas d'utilisation (résumé)}
\begin{longtable}{@{}p{3.5cm}p{11cm}@{}}
\toprule
\textbf{UC} & \textbf{Description} \\
\midrule
\endhead
UC1 & Consulter la liste, rechercher et filtrer des cartes. \\
UC2 & Ouvrir une carte, changer le pronom, écouter la prononciation. \\
UC3 & Imprimer une carte (grille par pronom). \\
UC4 & Créer une nouvelle carte avec validation. \\
UC5 & S'inscrire, se connecter, se déconnecter. \\
UC6 & (Admin) Personnaliser le thème et sauvegarder les réglages. \\
UC7 & (Admin) Consulter l'état du système et exécuter une action de maintenance. \\
UC8 & Consommer l'API \texttt{/api/verbs} depuis un client. \\
UC9 & Ajouter/retirer un favori et consulter la liste des favoris. \\
UC10 & Accéder au quiz et sélectionner des groupes. \\
UC11 & Dupliquer une carte existante pour en créer une nouvelle. \\
\bottomrule
\end{longtable}

\section{Critères d'acceptation (exemples)}
\begin{itemize}[leftmargin=*]
  \item La page \texttt{/cards} affiche une liste paginée et les filtres fonctionnent.
  \item La création refuse un infinitif du 1er groupe ne finissant pas par \texttt{er}.
  \item Deux verbes \texttt{manger} et \texttt{Manger} ne peuvent pas coexister.
  \item \texttt{/health} renvoie \texttt{status=ok} et HTTP 200 si DB+storage OK.
  \item Un utilisateur non admin reçoit HTTP 403 sur \texttt{/admin/theme} et \texttt{/admin/system}.
  \item La sauvegarde du thème met à jour l'apparence des cartes sans rechargement complet.
  \item Un utilisateur peut ajouter un favori puis le retrouver dans \texttt{/favorites}.
\end{itemize}

\section{Annexe : Références techniques (code)}
Pour valider la cohérence du cahier des charges avec l'implémentation, les références principales sont :
\begin{itemize}[leftmargin=*]
  \item Routes web : \texttt{routes/web.php}
  \item Routes API : \texttt{routes/api.php}
  \item Contrôleur cartes : \texttt{app/Http/Controllers/VerbCardController.php}
  \item Contrôleur thème : \texttt{app/Http/Controllers/ThemeController.php}
  \item Contrôleur système : \texttt{app/Http/Controllers/SystemController.php}
  \item Provider thème : \texttt{app/Providers/AppServiceProvider.php}
  \item JS global : \texttt{public/js/app.js}
\end{itemize}

\end{document}
